\chapter{The Graph Laplacian}

Every chapter so far has treated a graph as a combinatorial object: a set of vertices, a set of edges, and algorithms that crawl over them. This chapter performs a change of viewpoint that will carry us through the rest of the book. We attach to a graph a single symmetric matrix, the \emph{Laplacian} $\Lap$, and from that moment on the graph becomes a piece of linear algebra. Its eigenvalues will measure connectivity, its quadratic form will measure smoothness, its null space will count connected components, its pseudoinverse will compute electrical resistances and hitting times of random walks, and the linear system $\Lap x = b$ will become the computational engine behind spectral clustering, max-flow, and graph sparsification.

The Laplacian is the bridge between the discrete world of Parts I--V and the continuous world of spectra, electrical flows, and random walks that fills Parts VI and beyond. This chapter builds that bridge slowly and carefully. We define $\Lap = \Adjm_D - \Adjm$ three different ways, prove they coincide, and extract everything that follows from one master identity: $x^\top \Lap x$ equals the sum over edges of squared differences. That identity, more than any single theorem, is the thing to remember.

\begin{intuitionbox}
Think of the graph as a network of springs. Put a bead at each vertex and let bead $i$ sit at height $x_i$ on a vertical wire. Every edge is an ideal spring of stiffness $w_{ij}$ connecting the two beads it touches. The total elastic energy stored in the network, when the beads are pinned at heights $x = (x_1,\dots,x_n)$, is exactly
\[
  \tfrac12\, x^\top \Lap x \;=\; \tfrac12 \sum_{\{i,j\}\in E} w_{ij}\,(x_i - x_j)^2 .
\]
The Laplacian \emph{is} the energy of the spring network. Everything in this chapter --- smoothness, connectivity, the smallest eigenvalues --- is a statement about which height assignments store little energy. Flat assignments (all beads at the same height) store none; that is why the all-ones vector is in the null space.
\end{intuitionbox}

\begin{figure}[H]
\centering
\begin{tikzpicture}[thick]
  \coordinate (a) at (0,2.4);
  \coordinate (b) at (3,2.0);
  \coordinate (c) at (1.4,0.3);
  \draw[structurepurple,line width=1pt] (a) -- (b) node[midway,above,font=\scriptsize,black]{$w_{12}$};
  \draw[structurepurple,line width=1pt] (b) -- (c) node[midway,right,font=\scriptsize,black]{$w_{23}$};
  \draw[structurepurple,line width=1pt] (a) -- (c) node[midway,left,font=\scriptsize,black]{$w_{13}$};
  \foreach \p/\lab in {a/$x_1$,b/$x_2$,c/$x_3$}{
    \fill[nodeblue] (\p) circle (2.6pt);
    \node[font=\scriptsize,above right] at (\p) {\lab};
  }
  \draw[->,gray] (-0.8,-0.4) -- (-0.8,2.8) node[left,font=\scriptsize,black]{height};
\end{tikzpicture}
\caption{Three beads joined by springs of stiffness $w_{ij}$. Pinning the beads at heights $x_1,x_2,x_3$ stores elastic energy $\tfrac12\sum_{\{i,j\}}w_{ij}(x_i-x_j)^2 = \tfrac12 x^\top\Lap x$. Flattening every bead to a common height stores zero energy --- the all-ones null direction.}
\label{ch31:fig-springs}
\end{figure}

\chapterline

\section{Three definitions, one matrix}

Fix an undirected graph $\Gph=(V,E)$ with $n=\abs{V}$ vertices, labeled $1,\dots,n$, and $m=\abs{E}$ edges. For now take the graph unweighted and simple (no loops, no parallel edges); we generalize to weights in Section~\ref{ch31:weighted}. Recall two familiar matrices.

The \emph{adjacency matrix} $\Adjm \in \R^{n\times n}$ has $\Adjm_{ij}=1$ when $\{i,j\}\in E$ and $0$ otherwise. It is symmetric with zero diagonal. The \emph{degree matrix} $\Adjm_D \in \R^{n\times n}$ is diagonal, with $(\Adjm_D)_{ii} = \deg(i) = \sum_j \Adjm_{ij}$ on the diagonal and zeros off it.

\begin{definition}[Combinatorial Laplacian]\label{ch31:def-laplacian}
The \emph{Laplacian} of $\Gph$ is the $n\times n$ matrix
\[
  \Lap \;=\; \Adjm_D - \Adjm,
  \qquad
  \Lap_{ij} \;=\;
  \begin{cases}
    \deg(i) & i = j,\\
    -1 & \{i,j\}\in E,\\
    0 & \text{otherwise.}
  \end{cases}
\]
\end{definition}

Each diagonal entry holds a vertex's degree; each off-diagonal entry is $-1$ exactly when the two vertices are adjacent. Notice that every row sums to zero: the diagonal $\deg(i)$ is canceled by the $\deg(i)$ many $-1$'s in that row. We will see in a moment that this single observation is the algebraic heart of the matrix.

\begin{example}[A five-vertex graph]\label{ch31:ex-running}
We will trace the whole chapter on one little graph: a $4$-cycle $1\!-\!2\!-\!3\!-\!4\!-\!1$ with a pendant vertex $5$ hanging off vertex $3$.

\begin{figure}[H]
\centering
\begin{tikzpicture}[every node/.style={circle,draw,fill=nodeblue!15,minimum size=8mm,inner sep=1pt,font=\small},thick]
  \node (1) at (0,2) {1};
  \node (2) at (2,2) {2};
  \node (3) at (2,0) {3};
  \node (4) at (0,0) {4};
  \node (5) at (4,0) {5};
  \draw (1)--(2);
  \draw (2)--(3);
  \draw (3)--(4);
  \draw (4)--(1);
  \draw (3)--(5);
\end{tikzpicture}
\caption{Our running graph $\Gph$: a $4$-cycle with a pendant. Degrees are $\deg(1)=\deg(2)=\deg(4)=2$, $\deg(3)=3$, $\deg(5)=1$.}
\label{ch31:fig-running}
\end{figure}

Reading the edges $\{1,2\},\{2,3\},\{3,4\},\{1,4\},\{3,5\}$ straight into Definition~\ref{ch31:def-laplacian}:
\[
  \Adjm =
  \begin{pmatrix}
    0&1&0&1&0\\ 1&0&1&0&0\\ 0&1&0&1&1\\ 1&0&1&0&0\\ 0&0&1&0&0
  \end{pmatrix},
  \qquad
  \Lap =
  \begin{pmatrix}
    2&-1&0&-1&0\\ -1&2&-1&0&0\\ 0&-1&3&-1&-1\\ -1&0&-1&2&0\\ 0&0&-1&0&1
  \end{pmatrix}.
\]
Check the rows sum to zero. The diagonal $(2,2,3,2,1)$ is exactly the degree sequence. Keep this $\Lap$ in view; we will compute its quadratic form, its null space, and its eigenvalues by hand before the chapter ends.
\end{example}

\subsection{The incidence matrix and \texorpdfstring{$\Lap = \Adjm_B \Adjm_B^\top$}{L = B B\^T}}

The second definition is the one that explains \emph{why} $\Lap$ is positive semidefinite, and it is the one we lean on for proofs. To write it we must orient each edge --- pick, arbitrarily, a head and a tail. The orientation is a bookkeeping device; the final matrix $\Lap$ does not depend on it, as we will verify.

\begin{definition}[Incidence matrix]\label{ch31:def-incidence}
Orient each edge $e$ from a tail $t(e)$ to a head $h(e)$. The (signed, vertex--edge) \emph{incidence matrix} $\Adjm_B \in \R^{n\times m}$ has one row per vertex and one column per edge, with
\[
  (\Adjm_B)_{v,e} =
  \begin{cases}
    +1 & v = h(e),\\
    -1 & v = t(e),\\
    0 & \text{otherwise.}
  \end{cases}
\]
\end{definition}

Each column of $\Adjm_B$ has exactly one $+1$ and one $-1$: it is the vector $\mathbf{e}_{h(e)} - \mathbf{e}_{t(e)}$, where $\mathbf{e}_v$ is the standard basis vector. Geometrically the column of edge $e=(u\to v)$ is the ``discrete gradient direction'' along that edge.

\begin{theorem}[Incidence factorization]\label{ch31:thm-bbt}
For any orientation of the edges, $\Lap = \Adjm_B \Adjm_B^\top$.
\end{theorem}

\begin{whybox}
\textbf{Idea.} $\Adjm_B \Adjm_B^\top$ sums, over all edges, the outer product of each column with itself. A column for edge $\{u,v\}$ is $\mathbf{e}_u - \mathbf{e}_v$ (up to sign, which squares away), and $(\mathbf{e}_u-\mathbf{e}_v)(\mathbf{e}_u-\mathbf{e}_v)^\top$ places $+1$ at $(u,u)$ and $(v,v)$ and $-1$ at $(u,v)$ and $(v,u)$. Summing over edges builds up exactly the degree on the diagonal and a $-1$ for each adjacency. That is $\Lap$.
\end{whybox}

\begin{proof}
Write $\mathbf{b}_e \in \R^n$ for the column of $\Adjm_B$ indexed by edge $e$, so $\mathbf{b}_e = \mathbf{e}_{h(e)} - \mathbf{e}_{t(e)}$. Matrix multiplication gives
\[
  \Adjm_B \Adjm_B^\top \;=\; \sum_{e\in E} \mathbf{b}_e \mathbf{b}_e^\top .
\]
For a single edge $e$ with endpoints $u,v$ (in either orientation),
\[
  \mathbf{b}_e \mathbf{b}_e^\top
  = (\mathbf{e}_u - \mathbf{e}_v)(\mathbf{e}_u - \mathbf{e}_v)^\top
  = \mathbf{e}_u\mathbf{e}_u^\top + \mathbf{e}_v\mathbf{e}_v^\top - \mathbf{e}_u\mathbf{e}_v^\top - \mathbf{e}_v\mathbf{e}_u^\top .
\]
The sign of $\mathbf{b}_e$ does not matter: flipping the orientation negates $\mathbf{b}_e$, and $(-\mathbf{b}_e)(-\mathbf{b}_e)^\top = \mathbf{b}_e\mathbf{b}_e^\top$. Now sum over edges. The term $\mathbf{e}_u\mathbf{e}_u^\top$ contributes $+1$ to entry $(u,u)$, once for each edge incident to $u$; the total at $(u,u)$ is therefore $\deg(u)$. The cross terms $-\mathbf{e}_u\mathbf{e}_v^\top - \mathbf{e}_v\mathbf{e}_u^\top$ contribute $-1$ to entries $(u,v)$ and $(v,u)$, once for the single edge $\{u,v\}$ (the graph is simple). So $\big(\Adjm_B\Adjm_B^\top\big)_{ii} = \deg(i)$ and $\big(\Adjm_B\Adjm_B^\top\big)_{ij} = -1$ iff $\{i,j\}\in E$. That is precisely $\Lap$.
\end{proof}

\begin{example}[Incidence matrix of the running graph]\label{ch31:ex-incidence}
Orient each edge from its smaller to its larger label: $e_1=(1\to2)$, $e_2=(2\to3)$, $e_3=(3\to4)$, $e_4=(1\to4)$, $e_5=(3\to5)$. Then
\[
  \Adjm_B =
  \bordermatrix{
       & e_1 & e_2 & e_3 & e_4 & e_5 \cr
    1 & -1 & 0 & 0 & -1 & 0 \cr
    2 & +1 & -1 & 0 & 0 & 0 \cr
    3 & 0 & +1 & -1 & 0 & -1 \cr
    4 & 0 & 0 & +1 & +1 & 0 \cr
    5 & 0 & 0 & 0 & 0 & +1 \cr
  }.
\]
You should multiply $\Adjm_B\Adjm_B^\top$ and recover the $\Lap$ of Example~\ref{ch31:ex-running}; for instance row $3$ of $\Adjm_B$ has three nonzeros, so $(\Adjm_B\Adjm_B^\top)_{33}=3=\deg(3)$, and vertices $3,4$ share only edge $e_3$, where row $3$ holds $-1$ and row $4$ holds $+1$, so $(\Adjm_B\Adjm_B^\top)_{34}=(-1)(+1)=-1$ as required.
\end{example}

\subsection{The sum-over-edges form}

The third definition skips matrices entirely and writes down what the quadratic form $x^\top \Lap x$ \emph{does} to a vector. It is the most useful of the three.

\begin{theorem}[The quadratic form]\label{ch31:thm-quadform}
For every $x\in\R^n$,
\[
  x^\top \Lap x \;=\; \sum_{\{i,j\}\in E} (x_i - x_j)^2 \;\ge\; 0 .
\]
\end{theorem}

\begin{proof}
Use the factorization of Theorem~\ref{ch31:thm-bbt}: $x^\top \Lap x = x^\top \Adjm_B\Adjm_B^\top x = (\Adjm_B^\top x)^\top (\Adjm_B^\top x) = \norm{\Adjm_B^\top x}^2$. The vector $\Adjm_B^\top x \in \R^m$ has, in the coordinate of edge $e=(u\to v)$, the value $\mathbf{b}_e^\top x = x_v - x_u$. Hence
\[
  \norm{\Adjm_B^\top x}^2 = \sum_{e\in E} (x_{h(e)} - x_{t(e)})^2 = \sum_{\{i,j\}\in E}(x_i - x_j)^2 .
\]
The squares are orientation-independent and nonnegative, so $x^\top \Lap x \ge 0$.
\end{proof}

\begin{focusbox}
Three faces of the same matrix, used for three purposes:
\begin{itemize}
  \item $\Lap = \Adjm_D - \Adjm$ --- for \emph{building} $\Lap$ from a graph and reading off entries.
  \item $\Lap = \Adjm_B\Adjm_B^\top$ --- for \emph{proving} structural facts (PSD, rank, null space).
  \item $x^\top\Lap x = \sum_{\{i,j\}\in E}(x_i-x_j)^2$ --- for \emph{reasoning} about smoothness, cuts, and eigenvalues.
\end{itemize}
Memorize the third. Almost every estimate in spectral graph theory begins by writing $x^\top \Lap x$ as a sum of squared edge-differences and bounding it.
\end{focusbox}

\begin{historybox}
The word ``Laplacian'' is borrowed from continuous analysis. On a smooth domain the Laplace operator $\Delta f = \sum_i \partial^2 f/\partial x_i^2$ measures how much a function differs from its local average; its Dirichlet energy is $\int \norm{\nabla f}^2$. The graph Laplacian is the discrete analogue: $(\Lap x)_i = \sum_{j\sim i}(x_i - x_j)$ compares $x_i$ to its neighbors, and $x^\top \Lap x = \sum_{\{i,j\}}(x_i-x_j)^2$ is the discrete Dirichlet energy. The matrix itself appears in Kirchhoff's 1847 work on electrical networks --- his Matrix--Tree theorem (Chapter on spanning trees) counts spanning trees as a minor of $\Lap$ --- which is why it is sometimes called the \emph{Kirchhoff matrix}. The spectral perspective was developed by Fiedler (1973), whose ``algebraic connectivity'' $\lambda_2$ launched spectral graph theory, and pushed to its modern algorithmic frontier by Spielman, Teng, and others from the 2000s on.
\end{historybox}

\section{Smoothness: what \texorpdfstring{$x^\top \Lap x$}{x\^T L x} measures}

Read $x$ as a \emph{signal} on the vertices: a real number $x_i$ assigned to each vertex $i$. Theorem~\ref{ch31:thm-quadform} says $x^\top \Lap x$ penalizes disagreement across edges. A signal that is nearly constant on each tightly connected region has small Laplacian quadratic form; a signal that jumps wildly between neighbors has large form. That penalty is why the Laplacian appears in clustering, semi-supervised learning, mesh processing, and image segmentation.

\begin{intuitionbox}
The quantity $x^\top \Lap x$ is a \emph{roughness meter}. Feed it a constant signal and the needle reads zero --- there is no disagreement anywhere. Feed it a signal that respects the graph's community structure (high inside clusters, low between them, but changing only across the few inter-cluster edges) and the needle stays low. Feed it a checkerboard signal that flips sign across every edge and the needle pins to the maximum. The eigenvectors of $\Lap$ are, in order, the smoothest possible signals: the very smoothest is constant, the next is the gentlest non-constant variation the graph allows, and so on up to the roughest.
\end{intuitionbox}

Let us make ``smoothest non-constant signal'' precise, because that is the Fiedler vector that drives the next chapters. Because $\Lap$ is real symmetric, the spectral theorem (from the linear algebra companion) gives an orthonormal eigenbasis $v_1,\dots,v_n$ with real eigenvalues
\[
  0 = \lambda_1 \le \lambda_2 \le \cdots \le \lambda_n .
\]
The eigenvalues are nonnegative because $\Lap$ is PSD (Theorem~\ref{ch31:thm-quadform}), and $\lambda_1 = 0$ because the all-ones vector $\mathbf{1}$ kills the form. By the Courant--Fischer / Rayleigh-quotient characterization,
\[
  \lambda_2 \;=\; \min_{\substack{x \perp \mathbf{1} \\ x\neq 0}} \frac{x^\top \Lap x}{x^\top x}
  \;=\; \min_{\substack{x\perp\mathbf{1}\\ x\neq 0}} \frac{\sum_{\{i,j\}\in E}(x_i-x_j)^2}{\sum_i x_i^2}.
\]
So $\lambda_2$ is the least roughness achievable by any signal that is orthogonal to constants (i.e.\ has the constant component removed). Its minimizer $v_2$ is the \emph{Fiedler vector}: the smoothest way to assign real numbers to the vertices without making them all equal. We return to it at length in Chapter on spectral clustering; here we only need to see that the Laplacian's small eigenvalues encode the gentle, large-scale structure of the graph.

\begin{connectionbox}
This whole section is an instance of the linear-algebra book's chapter on \emph{symmetric PSD matrices and the spectral theorem}. Three imported facts do all the work: (1) a real symmetric matrix has a real orthonormal eigenbasis; (2) it is PSD iff $x^\top M x \ge 0$ for all $x$ iff all eigenvalues are $\ge 0$; (3) the Rayleigh quotient $x^\top M x / x^\top x$ ranges over $[\lambda_{\min},\lambda_{\max}]$, with the constrained minima given by Courant--Fischer. The graph Laplacian is the canonical example through which those abstractions become concrete and combinatorial. Eigenvalues become connectivity; eigenvectors become embeddings.
\end{connectionbox}

\begin{trapbox}
$\Lap$ is positive \emph{semi}definite, never positive definite, because $\mathbf{1}$ is always in its null space. So $\Lap$ is singular --- $\det \Lap = 0$ always --- and you cannot invert it. Algorithms that ``solve $\Lap x = b$'' really solve on the subspace $\mathbf{1}^\perp$, using the Moore--Penrose pseudoinverse $\Lap^{+}$, and they require the consistency condition $\mathbf{1}^\top b = 0$ (the entries of $b$ must sum to zero). Forgetting that the system is singular is the single most common Laplacian bug. We meet $\Lap^{+}$ properly in the electrical-networks chapter.
\end{trapbox}

\section{The null space and connected components}

The flat direction $\mathbf{1}$ is in the kernel of every Laplacian. The fact that pays off is that the \emph{entire} kernel is built from flat directions, one per connected component, and so the dimension of the null space counts the components. This is the first place where a purely algebraic quantity (a matrix rank) recovers a purely combinatorial one (the number of pieces).

\begin{theorem}[Kernel counts components]\label{ch31:thm-kernel}
Let $\Gph$ have connected components $C_1,\dots,C_k$. Then the null space of $\Lap$ has dimension exactly $k$, and is spanned by the indicator vectors $\mathbf{1}_{C_1},\dots,\mathbf{1}_{C_k}$ (the vector $\mathbf{1}_{C_\ell}$ is $1$ on the vertices of $C_\ell$ and $0$ elsewhere). Equivalently, $\rank \Lap = n - k$. In particular $\Gph$ is connected iff $\lambda_2 > 0$, and the multiplicity of the eigenvalue $0$ equals the number of components.
\end{theorem}

\begin{whybox}
\textbf{Idea.} $\Lap x = 0$ forces $x^\top \Lap x = 0$, and by the sum-of-squares form that forces $x_i = x_j$ across every edge. A signal that never changes along an edge is constant on each connected component but free to take a different value per component. So the kernel is exactly the $k$-dimensional space of ``locally constant'' signals.
\end{whybox}

\begin{proof}
We show $\ker \Lap = \operatorname{span}\set{\mathbf{1}_{C_1},\dots,\mathbf{1}_{C_k}}$.

\textbf{($\supseteq$)} Each $\mathbf{1}_{C_\ell}$ is constant on its component and zero elsewhere, so $(\mathbf{1}_{C_\ell})_i - (\mathbf{1}_{C_\ell})_j = 0$ for every edge $\{i,j\}$ (both endpoints lie in the same component, hence both get $1$, or both get $0$). By Theorem~\ref{ch31:thm-quadform}, $\mathbf{1}_{C_\ell}^\top \Lap \mathbf{1}_{C_\ell} = 0$. Because $\Lap$ is PSD, $x^\top \Lap x = 0 \iff \Lap x = 0$: indeed $0 = x^\top \Lap x = \norm{\Adjm_B^\top x}^2$ gives $\Adjm_B^\top x = 0$, hence $\Lap x = \Adjm_B(\Adjm_B^\top x) = 0$. So each $\mathbf{1}_{C_\ell}\in\ker\Lap$.

\textbf{($\subseteq$)} Suppose $\Lap x = 0$. Then $x^\top \Lap x = 0$, so $\sum_{\{i,j\}\in E}(x_i - x_j)^2 = 0$, which forces $x_i = x_j$ for every edge $\{i,j\}$. Pick any vertex and follow paths: within a connected component every pair of vertices is joined by a path, and $x$ is unchanged along each path's edges, so $x$ is constant on the component. Therefore $x = \sum_\ell c_\ell \mathbf{1}_{C_\ell}$ for some scalars $c_\ell$ (the value $x$ takes on $C_\ell$). Hence $x\in\operatorname{span}\set{\mathbf{1}_{C_\ell}}$.

The indicators $\mathbf{1}_{C_1},\dots,\mathbf{1}_{C_k}$ have disjoint supports, so they are linearly independent; thus $\dim\ker\Lap = k$. By rank--nullity $\rank\Lap = n-k$. The eigenvalue $0$ has multiplicity $\dim\ker\Lap = k$; ordering eigenvalues $0=\lambda_1\le\cdots\le\lambda_n$, connectivity ($k=1$) is equivalent to $\lambda_2 > 0$.
\end{proof}

\begin{labbox}
Run this on Example~\ref{ch31:ex-running}. The graph is connected, so $k=1$ and the kernel should be one-dimensional, spanned by $\mathbf{1}=(1,1,1,1,1)^\top$. Check: each row of $\Lap$ sums to zero, which is exactly the statement $\Lap\mathbf{1}=0$. Now imagine deleting the single edge $\{3,5\}$, which strands vertex $5$ on its own. With two components the kernel jumps to dimension $2$, spanned by $\mathbf{1}_{\{1,2,3,4\}}$ and $\mathbf{1}_{\{5\}}$, and $\rank\Lap$ falls to $n-2=3$. The rank of $\Lap$ drops by exactly one for each new component.
\end{labbox}

\subsection{An algorithm: counting components by Laplacian rank}

Theorem~\ref{ch31:thm-kernel} turns ``how many connected components?'' into ``what is the rank of $\Lap$?''. For a real graph you would just run BFS/DFS in $\Oh(n+m)$ --- nobody computes a rank to count components. But the reduction is conceptually important (it is how spectral methods ``see'' connectivity), and the same Gaussian elimination is the skeleton of the Laplacian solvers to come. Here is the explicit procedure.

\begin{algorithm}[H]
\caption{\textsc{ComponentsByRank}$(\Gph)$ --- count components via $\dim\ker\Lap$}
\label{ch31:alg-rank}
\DontPrintSemicolon
\KwIn{Graph $\Gph=(V,E)$ on $n$ vertices.}
\KwOut{Number of connected components $k$.}
Build $\Lap = \Adjm_D - \Adjm \in \R^{n\times n}$.\;
$r \gets 0$\tcp*{rank accumulator}
\For{$c \gets 1$ \KwTo $n$ (sweep columns)}{
  Find a pivot row $p$ with $\abs{\Lap_{p,c}}$ maximal among rows not yet used\;
  \If{$\abs{\Lap_{p,c}} > \tau$ \tcp*{numerical zero tolerance $\tau$}}{
    Mark row $p$ used; $r \gets r+1$\;
    \For{every unused row $q$}{
      $\Lap_{q,\bullet} \gets \Lap_{q,\bullet} - \dfrac{\Lap_{q,c}}{\Lap_{p,c}}\,\Lap_{p,\bullet}$ \tcp*{eliminate column $c$ below the pivot}
    }
  }
}
\Return $k \gets n - r$\;
\end{algorithm}

\begin{whybox}
The algorithm is plain Gaussian elimination with partial pivoting; the number of pivots it finds is the rank $r$ of $\Lap$. By Theorem~\ref{ch31:thm-kernel}, $\rank\Lap = n-k$, so the missing $n-r$ pivots are exactly the $k$ components. Each ``rank-deficient'' column --- a column with no pivot --- corresponds to one degree of freedom in the null space, i.e.\ one free component value. Correctness is therefore immediate from the theorem; the only subtlety is numerical (choosing $\tau$ so that genuine zeros are not mistaken for tiny nonzeros), which is why partial pivoting matters.
\end{whybox}

\begin{complexitybox}
Building $\Lap$ costs $\Oh(n^2)$ to allocate plus $\Oh(m)$ to fill. Dense Gaussian elimination on an $n\times n$ matrix is $\Oh(n^3)$ time and $\Oh(n^2)$ space --- dominant. So \textsc{ComponentsByRank} runs in $\Th(n^3)$. This is dramatically worse than the $\Oh(n+m)$ of BFS, which is the point: \emph{never} use rank to count components in practice. The value of the reduction is conceptual and it foreshadows fast Laplacian system solvers, which exploit sparsity to beat $\Oh(n^3)$ by orders of magnitude.
\end{complexitybox}

\begin{labbox}
Trace \textsc{ComponentsByRank} on the running graph's $\Lap$ to watch the rank come out at $4 = n-1$. We sweep columns left to right; after each pivot we clear the entries below it. Starting from
\[
  \Lap =
  \begin{pmatrix}
    2&-1&0&-1&0\\ -1&2&-1&0&0\\ 0&-1&3&-1&-1\\ -1&0&-1&2&0\\ 0&0&-1&0&1
  \end{pmatrix},
\]
column $1$ pivots on the $2$ in row $1$; eliminating below gives rows $2,4$ updated and a $2\times2$ Schur block forming. Carrying the elimination through columns $1$--$4$ (each finds a pivot above tolerance $\tau$) yields an upper-triangular shape with exactly four nonzero pivots and a final all-zero row:
\[
  \xrightarrow{\text{eliminate}}
  \begin{pmatrix}
    2&-1&0&-1&0\\ 0&\tfrac32&-1&-\tfrac12&0\\ 0&0&\tfrac{7}{3}&-\tfrac43&-1\\ 0&0&0&\tfrac47&-\tfrac47\\ 0&0&0&0&0
  \end{pmatrix}.
\]
Four pivots ($2,\tfrac32,\tfrac73,\tfrac47$), so $r=4$ and $k = n-r = 5-4 = 1$: one component, as expected. The lone zero pivot in the last row is the all-ones null direction making itself felt --- the single ``free'' degree of freedom the theorem promised. Add a second component (snip vertex $5$ free) and a second zero pivot appears.
\end{labbox}

\section{Weighted Laplacians}\label{ch31:weighted}

Real networks carry weights: edge $\{i,j\}$ has a conductance $w_{ij} > 0$ (think spring stiffness, pipe capacity, or similarity score). Everything above generalizes with the weights tucked into the right places, and not a single proof changes in spirit.

\begin{definition}[Weighted Laplacian]\label{ch31:def-weighted}
Let $w:E\to\R_{>0}$ assign a positive weight to each edge. The \emph{weighted degree} of $i$ is $d_i = \sum_{j:\{i,j\}\in E} w_{ij}$. The \emph{weighted adjacency} $\Adjm$ has $\Adjm_{ij} = w_{ij}$ for edges and $0$ otherwise; the \emph{weighted degree matrix} $\Adjm_D$ is diagonal with $d_i$. The weighted Laplacian is
\[
  \Lap = \Adjm_D - \Adjm,
  \qquad
  \Lap_{ij} =
  \begin{cases}
    d_i & i=j,\\
    -w_{ij} & \{i,j\}\in E,\\
    0 & \text{otherwise.}
  \end{cases}
\]
\end{definition}

The factorization survives with a diagonal weight matrix in the middle. Let $\Adjm_W = \operatorname{diag}(w_e)_{e\in E} \in \R^{m\times m}$ hold the edge weights and let $\Adjm_B$ be the (unweighted) signed incidence matrix.

\begin{theorem}[Weighted factorization and quadratic form]\label{ch31:thm-weighted}
With $\Adjm_B,\Adjm_W$ as above,
\[
  \Lap = \Adjm_B\, \Adjm_W\, \Adjm_B^\top
  = \sum_{e\in E} w_e\, \mathbf{b}_e\mathbf{b}_e^\top,
  \qquad
  x^\top \Lap x = \sum_{\{i,j\}\in E} w_{ij}\,(x_i - x_j)^2 \ge 0 .
\]
Consequently $\Lap$ is PSD, its kernel is spanned by the component indicators (Theorem~\ref{ch31:thm-kernel} holds verbatim, since $w_{ij}>0$ makes $(x_i-x_j)^2=0$ exactly when adjacent), and the spring-energy reading $\tfrac12 x^\top\Lap x$ is the elastic energy with stiffnesses $w_{ij}$.
\end{theorem}

\begin{proof}
$\Adjm_B\Adjm_W\Adjm_B^\top = \sum_e w_e\,\mathbf{b}_e\mathbf{b}_e^\top$ because $\Adjm_W$ is diagonal: the $e$-th column of $\Adjm_B$ gets scaled by $w_e$ before the outer product. Each $w_e\,\mathbf{b}_e\mathbf{b}_e^\top$ places $+w_e$ on the two diagonal entries of its endpoints and $-w_e$ on the two off-diagonal entries; summing yields $d_i$ on the diagonal and $-w_{ij}$ off it, matching Definition~\ref{ch31:def-weighted}. The quadratic form follows as in Theorem~\ref{ch31:thm-quadform}: writing $\Adjm_W^{1/2}\Adjm_B^\top x$ and taking its squared norm gives $\sum_e w_e (x_{h(e)}-x_{t(e)})^2$. Since each $w_e>0$, the form vanishes iff every adjacent pair agrees, and the kernel argument of Theorem~\ref{ch31:thm-kernel} carries over unchanged.
\end{proof}

\begin{example}[A weighted twist on the running graph]\label{ch31:ex-weighted}
Keep the graph of Figure~\ref{ch31:fig-running} but weight the cycle edges by $2$ and the pendant edge $\{3,5\}$ by $5$:
\[
  w_{12}=w_{23}=w_{34}=w_{14}=2,\quad w_{35}=5.
\]
Then $d_1=d_2=d_4 = 2+2 = 4$, $d_3 = 2+2+5 = 9$, $d_5 = 5$, and
\[
  \Lap =
  \begin{pmatrix}
    4 & -2 & 0 & -2 & 0\\
    -2 & 4 & -2 & 0 & 0\\
    0 & -2 & 9 & -2 & -5\\
    -2 & 0 & -2 & 4 & 0\\
    0 & 0 & -5 & 0 & 5
  \end{pmatrix}.
\]
Rows still sum to zero, so $\Lap\mathbf{1}=0$. Take the signal $x = (1,1,1,1,0)^\top$ that flags the pendant vertex. Then
\[
  x^\top \Lap x = \sum_{\{i,j\}} w_{ij}(x_i-x_j)^2 = w_{35}(1-0)^2 = 5,
\]
because the only edge whose endpoints disagree is $\{3,5\}$. Every cycle edge contributes $0$. The form reads off the cost of separating vertex $5$ --- a preview of the cut interpretation below.
\end{example}

\subsection{Cuts as quadratic forms}

A special signal makes the cut connection exact. Given a vertex set $S\subseteq V$, let $\mathbf{1}_S \in \set{0,1}^n$ be its indicator. Plug it in:
\[
  \mathbf{1}_S^\top \Lap\, \mathbf{1}_S
  = \sum_{\{i,j\}\in E} w_{ij}\big((\mathbf{1}_S)_i - (\mathbf{1}_S)_j\big)^2 .
\]
An edge contributes $w_{ij}$ exactly when one endpoint is in $S$ and the other is not --- that is, exactly when the edge \emph{crosses} the cut $(S,\bar S)$. So
\[
  \boxed{\;\mathbf{1}_S^\top \Lap\, \mathbf{1}_S = w(S,\bar S)\;}
\]
the Laplacian quadratic form evaluated at a $0/1$ indicator equals the weight of the corresponding cut. This is the seed of spectral cut bounds: relaxing $\mathbf{1}_S\in\set{0,1}^n$ to a real vector $x\perp\mathbf{1}$ turns the combinatorial minimum cut into the eigenvalue $\lambda_2$, and the gap between them is governed by Cheeger's inequality (next chapters). Min-cut, eigenvalue, electrical flow: all three are readings of the same quadratic form.

\begin{applicationbox}
This identity is why the Laplacian sits under so much applied machinery.
\begin{itemize}
  \item \textbf{Spectral clustering and image segmentation.} Minimizing a normalized version of $\mathbf{1}_S^\top\Lap\mathbf{1}_S$ over balanced $S$ is NP-hard; relaxing the indicator to the Fiedler vector $v_2$ gives a polynomial-time cut that Cheeger's inequality certifies as near-optimal. Shi--Malik's normalized cuts segment images exactly this way.
  \item \textbf{Semi-supervised learning.} Given a few labeled vertices, the ``harmonic'' extension minimizes $x^\top\Lap x$ subject to matching the labels --- the smoothest signal consistent with what you know. This is a single (constrained) Laplacian solve.
  \item \textbf{Graph drawing.} Plotting each vertex at coordinates $(v_2(i), v_3(i))$ from the second and third eigenvectors produces the classic ``spectral layout'' that lays out meshes and social graphs with edges short and clusters visible.
  \item \textbf{Mesh processing.} The cotangent Laplacian (a weighted Laplacian with geometry-derived $w_{ij}$) drives smoothing, parameterization, and the discrete heat equation in computer graphics.
\end{itemize}
\end{applicationbox}

\section{Eigenvalues of small graphs, by hand}

Abstract bounds become real once you compute a spectrum. We do two by hand and read the running example off them.

\begin{example}[The path and the cycle]\label{ch31:ex-path-cycle}
For the cycle $C_n$ every vertex has degree $2$, so $\Lap = 2I - \Adjm$ where $\Adjm$ is the cyclic shift plus its inverse. The eigenvectors are the discrete Fourier modes $v^{(\omega)}_j = e^{2\pi i \omega j/n}$ for $\omega=0,1,\dots,n-1$, and a direct computation gives
\[
  \lambda_\omega = 2 - 2\cos\!\Big(\frac{2\pi\omega}{n}\Big), \qquad \omega = 0,\dots,n-1 .
\]
The smallest is $\lambda_0 = 0$ (the constant mode), and the second smallest is $\lambda_2 = 2 - 2\cos(2\pi/n) \approx (2\pi/n)^2$ for large $n$ --- the algebraic connectivity of a long cycle shrinks like $1/n^2$, matching the intuition that a big ring is only weakly connected end-to-end. For the path $P_n$ the eigenvalues are $2 - 2\cos(\pi\omega/n)$, $\omega=0,\dots,n-1$, via cosine modes that respect the free (zero-flux) boundary.
\end{example}

\begin{example}[The complete graph and a star]\label{ch31:ex-kn-star}
For $K_n$ every off-diagonal entry is $-1$ and every diagonal is $n-1$, so $\Lap = nI - \mathbf{1}\mathbf{1}^\top$. The vector $\mathbf{1}$ gives eigenvalue $0$; any vector orthogonal to $\mathbf{1}$ gives $\Lap x = nx$. Hence the spectrum is $0$ once and $n$ with multiplicity $n-1$. The huge spectral gap $\lambda_2 = n$ says $K_n$ is maximally well connected --- you cannot cut it cheaply. The star $K_{1,n-1}$ has spectrum $0$, then $1$ with multiplicity $n-2$, then $n$ once; the lone large eigenvalue is the center's high degree.
\end{example}

\begin{labbox}
Now pin down the running graph (Example~\ref{ch31:ex-running}). It is connected, so $\lambda_1 = 0 < \lambda_2$. Two quick sanity checks bracket the spectrum without solving the characteristic polynomial:
\begin{itemize}
  \item \textbf{Trace} $=\sum_i \lambda_i = \sum_i \deg(i) = 2+2+3+2+1 = 10 = 2m$. (The trace of $\Lap$ is $\sum_i \deg i = 2m$, always.)
  \item \textbf{Largest eigenvalue} $\lambda_n \le 2\,\Delta = 6$ where $\Delta=\max_i\deg i = 3$ (a standard Gershgorin bound: every Gershgorin disc is centered at $\deg i$ with radius $\deg i$, so $\lambda \le 2\deg i$). Also $\lambda_n \ge \Delta + 1 = 4$ for any graph with a vertex of degree $\Delta$.
\end{itemize}
Solving numerically, the five eigenvalues are approximately $\{0,\ 0.83,\ 2.00,\ 2.69,\ 4.48\}$. They sum to $10$ (trace check passes), the gap $\lambda_2\approx0.83$ is small because the pendant vertex $5$ is loosely attached, and the Fiedler vector $v_2$ has its extreme entry at vertex $5$ --- the natural ``odd one out'' that the eigenvector pulls farthest from the cycle. Confirm the sum yourself; finding $v_2$ by hand is Exercise~\ref{ch31:ex-fiedler-hand}.
\end{labbox}

\begin{figure}[H]
\centering
\begin{tikzpicture}[scale=1.0]
  \draw[->,thick] (-0.3,0) -- (5.2,0) node[right]{$\lambda$};
  \foreach \x/\lab in {0/0, 0.83/{0.83}, 2/{2.0}, 2.69/{2.69}, 4.48/{4.48}}{
    \draw[thick,nodeblue] (\x,-0.12) -- (\x,0.12);
    \node[below,font=\scriptsize] at (\x,-0.16) {\lab};
  }
  \node[above,font=\footnotesize] at (0,0.2) {$\lambda_1$};
  \node[above,font=\footnotesize] at (0.83,0.2) {$\lambda_2$};
  \node[above,font=\footnotesize] at (4.48,0.2) {$\lambda_5$};
  \draw[<->,structurepurple,thick] (0,0.55) -- (0.83,0.55);
  \node[above,font=\scriptsize,structurepurple] at (0.42,0.55) {spectral gap};
\end{tikzpicture}
\caption{The Laplacian spectrum of the running graph on the real line. The gap from $\lambda_1=0$ to $\lambda_2\approx0.83$ is the algebraic connectivity; the small gap reflects the weakly attached pendant.}
\label{ch31:fig-spectrum}
\end{figure}

\section{Basic spectral facts to carry forward}

We collect the elementary identities that every later chapter assumes. They all fall out of the three definitions.

\begin{proposition}[Spectral toolbox]\label{ch31:prop-toolbox}
Let $\Lap$ be the (weighted) Laplacian of $\Gph$ with eigenvalues $0=\lambda_1\le\cdots\le\lambda_n$. Then:
\begin{enumerate}
  \item $\sum_{i} \lambda_i = \operatorname{tr}\Lap = \sum_i d_i = 2\sum_e w_e$ (twice the total edge weight).
  \item $\lambda_n \le 2\max_i d_i$, and $\lambda_n \ge \max_i d_i + 1$ for unweighted graphs with at least one edge.
  \item $\lambda_2 > 0 \iff \Gph$ is connected; the multiplicity of $0$ is the number of components.
  \item (Monotonicity) Adding an edge or increasing a weight can only increase each $\lambda_i$: if $\Lap' = \Lap + w\,\mathbf{b}_e\mathbf{b}_e^\top$ with $w\ge0$ then $\lambda_i(\Lap') \ge \lambda_i(\Lap)$ for all $i$.
  \item (Kirchhoff / Matrix--Tree) The number of spanning trees of a connected $\Gph$ equals $\tfrac1n \prod_{i=2}^{n}\lambda_i$, equivalently any cofactor of $\Lap$.
\end{enumerate}
\end{proposition}

\begin{proof}[Proof sketches]
(1) The trace of $\Lap$ is $\sum_i d_i$ by Definition~\ref{ch31:def-weighted}, and the trace equals the eigenvalue sum for any symmetric matrix. Each edge contributes its weight to two diagonal entries, so $\sum_i d_i = 2\sum_e w_e$.
(2) By Gershgorin, every eigenvalue lies in some disc centered at $d_i$ with radius $\sum_{j\neq i}\abs{\Lap_{ij}} = d_i$, so $\lambda \le 2d_i \le 2\max_i d_i$. The lower bound follows from interlacing with the star subgraph at a max-degree vertex.
(3) Theorem~\ref{ch31:thm-kernel}.
(4) Adding $w\,\mathbf{b}_e\mathbf{b}_e^\top$ adds a PSD rank-one matrix; by the Courant--Fischer min--max characterization, adding a PSD matrix raises every $\lambda_i$ (Weyl's monotonicity). Concretely the Rayleigh quotient $x^\top\Lap x/x^\top x$ increases pointwise.
(5) This is Kirchhoff's theorem; its proof via the Cauchy--Binet expansion of a cofactor of $\Adjm_B\Adjm_B^\top$ appears in the spanning-trees material. The product form follows because the nonzero eigenvalues are $\lambda_2,\dots,\lambda_n$ and a cofactor equals $\tfrac1n\prod_{i\ge2}\lambda_i$.
\end{proof}

\begin{connectionbox}
Item (5) is a thread tying this chapter to three others. The \emph{spanning-trees} chapter proves Matrix--Tree by Cauchy--Binet on $\Adjm_B\Adjm_B^\top$; the \emph{electrical networks} chapter reads $\Lap^{+}$ as effective resistance and shows that effective resistance equals the fraction of spanning trees through an edge; and \emph{Karger's} random-contraction analysis (global min-cut chapter) is, through resistance, secretly a statement about $\Lap$. The Laplacian is the common substrate. The linear-algebra companion supplies the spectral theorem, Courant--Fischer, Weyl monotonicity, and the pseudoinverse that all of this rests on.
\end{connectionbox}

\subsection{The normalized Laplacian, briefly}

Two variants normalize away the degrees, which matters when degrees vary wildly (web graphs, social networks). With $\Adjm_D$ the degree matrix,
\[
  \Lap_{\mathrm{sym}} = \Adjm_D^{-1/2}\,\Lap\,\Adjm_D^{-1/2} = I - \Adjm_D^{-1/2}\Adjm\,\Adjm_D^{-1/2},
  \qquad
  \Lap_{\mathrm{rw}} = \Adjm_D^{-1}\Lap = I - \Adjm_D^{-1}\Adjm .
\]
$\Lap_{\mathrm{sym}}$ is symmetric PSD with eigenvalues in $[0,2]$; $\Lap_{\mathrm{rw}}$ is the generator of the lazy random walk and shares the same eigenvalues (they are similar matrices). Its quadratic form is
\[
  x^\top \Lap_{\mathrm{sym}}\, x = \sum_{\{i,j\}\in E} w_{ij}\Big(\frac{x_i}{\sqrt{d_i}} - \frac{x_j}{\sqrt{d_j}}\Big)^2,
\]
so a signal is ``smooth'' relative to degree. The normalized spectrum is what Cheeger's inequality and the random-walk mixing bounds actually use, so we flag it here and develop it fully in the spectral-clustering and random-walk chapters.

\begin{trapbox}
$\Lap_{\mathrm{rw}} = \Adjm_D^{-1}\Lap$ is \emph{not symmetric}, so you must not feed it to a symmetric eigensolver or assume orthogonal eigenvectors. It is similar to the symmetric $\Lap_{\mathrm{sym}}$ via $\Lap_{\mathrm{rw}} = \Adjm_D^{-1/2}\Lap_{\mathrm{sym}}\Adjm_D^{1/2}$, so they share eigenvalues, but their eigenvectors differ by the factor $\Adjm_D^{\pm1/2}$. Compute with $\Lap_{\mathrm{sym}}$, then transform. Also: $\Adjm_D^{-1/2}$ needs every degree positive, so isolated vertices (degree $0$) must be removed first.
\end{trapbox}

\section{A worked solve: harmonic extension}

To make the algebra do something, we solve a small Laplacian system by hand: the harmonic-extension problem behind semi-supervised learning and electrical potentials. Fix the values of $x$ on a set $B\subset V$ of ``boundary'' (labeled) vertices and choose the rest to minimize $x^\top\Lap x$. The minimizer is \emph{harmonic} on the interior $I = V\setminus B$: each interior value equals the (weighted) average of its neighbors.

\begin{theorem}[Harmonic extension]\label{ch31:thm-harmonic}
Partition $\Lap$ in block form by interior $I$ and boundary $B$:
\[
  \Lap = \begin{pmatrix} \Lap_{II} & \Lap_{IB}\\ \Lap_{BI} & \Lap_{BB}\end{pmatrix},
  \qquad
  x = \begin{pmatrix} x_I \\ x_B \end{pmatrix}.
\]
With $x_B$ fixed, $x^\top\Lap x$ is minimized by the unique solution of $\Lap_{II}\,x_I = -\Lap_{IB}\,x_B$, provided every component of $\Gph$ touches $B$ (so $\Lap_{II}$ is invertible). At the optimum $(\Lap x)_i = 0$ for all $i\in I$: each interior vertex satisfies $d_i x_i = \sum_{j\sim i} w_{ij} x_j$, the weighted-average (harmonic) condition.
\end{theorem}

\begin{proof}
Expand $f(x_I) = x^\top\Lap x = x_I^\top\Lap_{II}x_I + 2 x_I^\top\Lap_{IB}x_B + x_B^\top\Lap_{BB}x_B$ (using symmetry $\Lap_{BI}=\Lap_{IB}^\top$). This is a convex quadratic in $x_I$ since $\Lap_{II}\succeq0$. Setting the gradient to zero, $\nabla_{x_I} f = 2\Lap_{II}x_I + 2\Lap_{IB}x_B = 0$, i.e.\ $\Lap_{II}x_I = -\Lap_{IB}x_B$. If every component meets $B$, then $\Lap_{II}$ is a principal submatrix of $\Lap$ with no all-zero-row support left in the kernel, hence positive \emph{definite} and invertible, giving a unique minimizer. The optimality condition $(\Lap x)_i = 0$ for $i\in I$ unpacks to $d_i x_i - \sum_{j\sim i}w_{ij}x_j = 0$.
\end{proof}

\begin{example}[Solving on the running graph]\label{ch31:ex-harmonic}
Use the unweighted $\Lap$ of Example~\ref{ch31:ex-running}. Pin the boundary $B=\set{1,5}$ with $x_1 = 0$ and $x_5 = 1$ (think: vertex $1$ grounded, vertex $5$ at unit potential), and solve for the interior $I=\set{2,3,4}$. The interior block and coupling are
\[
  \Lap_{II} = \begin{pmatrix} 2&-1&0\\ -1&3&-1\\ 0&-1&2\end{pmatrix}\ (\text{rows }2,3,4),
  \qquad
  -\Lap_{IB}x_B = -\begin{pmatrix} -1 & 0\\ 0 & -1\\ -1 & 0\end{pmatrix}\begin{pmatrix}0\\1\end{pmatrix}
  = \begin{pmatrix}0\\1\\0\end{pmatrix}.
\]
(The columns of $\Lap_{IB}$ are vertices $1$ and $5$; vertex $3$ couples to vertex $5$ with $-1$.) Solving $\Lap_{II}x_I = (0,1,0)^\top$:
\[
  2x_2 = x_3,\quad 3x_3 = x_2 + x_4 + 1,\quad 2x_4 = x_3 .
\]
By symmetry $x_2 = x_4$, so $x_3 = 2x_2$ and $3(2x_2) = 2x_2 + 1 \Rightarrow 4x_2 = 1 \Rightarrow x_2 = x_4 = \tfrac14$, $x_3 = \tfrac12$. The harmonic signal is $x = (0,\ \tfrac14,\ \tfrac12,\ \tfrac14,\ 1)$. Sanity check at the interior: vertex $3$ has $d_3=3$ and neighbors $2,4,5$ with values $\tfrac14,\tfrac14,1$, summing to $\tfrac32 = 3\cdot\tfrac12 = d_3 x_3$. Harmonic. This is exactly the electrical potential when a unit voltage is applied across vertices $1$ and $5$; the energy $x^\top\Lap x$ is the dissipated power, and its reciprocal is the effective conductance --- the bridge to the electrical-networks chapter.
\end{example}

\begin{figure}[H]
\centering
\begin{tikzpicture}[every node/.style={circle,draw,minimum size=9mm,inner sep=1pt,font=\small},thick]
  \node[fill=nodeblue!10] (1) at (0,2) {$0$};
  \node[fill=nodeblue!25] (2) at (2,2) {$\tfrac14$};
  \node[fill=nodeblue!45] (3) at (2,0) {$\tfrac12$};
  \node[fill=nodeblue!25] (4) at (0,0) {$\tfrac14$};
  \node[fill=nodeblue!80,text=white] (5) at (4,0) {$1$};
  \draw (1)--(2); \draw (2)--(3); \draw (3)--(4); \draw (4)--(1); \draw (3)--(5);
  \node[draw=none,below=2mm of 1,font=\scriptsize] {pinned};
  \node[draw=none,below=2mm of 5,font=\scriptsize] {pinned};
\end{tikzpicture}
\caption{The harmonic extension of Example~\ref{ch31:ex-harmonic}. Boundary vertices $1$ and $5$ are pinned at $0$ and $1$ (darker shading is higher potential); each interior value is the average of its neighbors. Read the picture as the steady voltage when a one-volt battery is wired across vertices $1$ and $5$ through unit resistors.}
\label{ch31:fig-harmonic}
\end{figure}

\begin{algorithm}[H]
\caption{\textsc{HarmonicExtend}$(\Lap, B, x_B)$ --- smoothest interior given boundary}
\label{ch31:alg-harmonic}
\DontPrintSemicolon
\KwIn{Laplacian $\Lap$; boundary set $B$ with fixed values $x_B$; interior $I = V\setminus B$.}
\KwOut{Signal $x$ minimizing $x^\top\Lap x$ subject to $x|_B = x_B$.}
Form blocks $\Lap_{II}$ (rows/cols $I$) and $\Lap_{IB}$ (rows $I$, cols $B$).\;
Assert every component of $\Gph$ meets $B$ \tcp*{else $\Lap_{II}$ singular}
Solve the SPD system $\Lap_{II}\, x_I = -\Lap_{IB}\, x_B$ \tcp*{Cholesky, CG, or a fast Laplacian solver}
$x \gets$ concatenate $x_I$ and $x_B$ in vertex order\;
\Return $x$\;
\end{algorithm}

\begin{complexitybox}
The cost is one solve of an SPD system of size $\abs{I}\le n$. Dense Cholesky is $\Oh(\abs{I}^3)$. Conjugate gradient runs in $\Oh(m\sqrt{\kappa})$ where $\kappa$ is the condition number of $\Lap_{II}$, which can be large. The headline result of modern spectral graph theory --- Spielman--Teng and successors --- gives \emph{nearly-linear} solvers running in $\Oh(m\log^{c}n\log(1/\eps))$ time for $\Lap x = b$ to accuracy $\eps$, by preconditioning with a spectral sparsifier or a low-stretch spanning tree. Those solvers are the computational payoff of this whole chapter and are developed in the spectral-methods part.
\end{complexitybox}

\section*{Exercises}
\addcontentsline{toc}{section}{Exercises}

\begin{exercise}[Build it]\label{ch31:ex-build}
Write down the Laplacian of the $3\times3$ grid graph (nine vertices, edges between orthogonally adjacent cells). Verify by hand that every row sums to zero and that the diagonal is the degree sequence $(2,3,2,3,4,3,2,3,2)$. Confirm $\operatorname{tr}\Lap = 2m$ by counting edges.
\end{exercise}

\begin{exercise}[Orientation independence]\label{ch31:ex-orient}
Prove directly from Definition~\ref{ch31:def-incidence} that $\Adjm_B\Adjm_B^\top$ does not depend on the chosen edge orientations, without invoking Theorem~\ref{ch31:thm-bbt}. (Hint: flipping one edge negates one column of $\Adjm_B$; show the outer product is unchanged.)
\end{exercise}

\begin{exercise}[Disconnect and watch the rank]\label{ch31:ex-rank-drop}
Take the running graph (Example~\ref{ch31:ex-running}) and delete edges until it has exactly three components. Predict $\dim\ker\Lap$ and $\rank\Lap$ before computing, then verify by row-reducing the new $\Lap$. Identify the three indicator vectors spanning the kernel.
\end{exercise}

\begin{exercise}[Cut = quadratic form]\label{ch31:ex-cut-form}
For the weighted graph of Example~\ref{ch31:ex-weighted}, compute $\mathbf{1}_S^\top\Lap\,\mathbf{1}_S$ for $S = \set{1,2}$, $S=\set{5}$, and $S=\set{1,2,4}$. Check each against the total weight of edges crossing the cut, and find the minimum-weight cut by trying all $2^4-1$ nontrivial vertex sets containing vertex $1$.
\end{exercise}

\begin{exercise}[Cycle spectrum]\label{ch31:ex-cycle}
Verify the cycle eigenvalue formula $\lambda_\omega = 2-2\cos(2\pi\omega/n)$ of Example~\ref{ch31:ex-path-cycle} for $C_4$ by computing the $4\times4$ Laplacian's characteristic polynomial. Confirm the spectrum is $\set{0,2,2,4}$, and identify which eigenvector is the Fiedler vector. Why is $\lambda_2$ degenerate (multiplicity $2$) here?
\end{exercise}

\begin{exercise}[Fiedler vector by hand]\label{ch31:ex-fiedler-hand}
For the running graph, the claim was $\lambda_2 \approx 0.83$ with the Fiedler vector taking its extreme value at vertex $5$. Set up the eigenproblem $\Lap v = \lambda v$ with $v\perp\mathbf{1}$, exploit the symmetry $x_2=x_4$, and reduce to a small polynomial. Solve for $\lambda_2$ to two decimals and confirm the sign pattern of $v_2$ separates the pendant from the cycle.
\end{exercise}

\begin{exercise}[PSD without incidence]\label{ch31:ex-psd-direct}
Prove $\Lap\succeq0$ directly from Definition~\ref{ch31:def-laplacian} (degree minus adjacency), by completing the square on $x^\top\Lap x = \sum_i \deg(i)x_i^2 - 2\sum_{\{i,j\}}x_ix_j$. (You should recover $\sum_{\{i,j\}}(x_i-x_j)^2$.)
\end{exercise}

\begin{exercise}[Weighted monotonicity]\label{ch31:ex-monotone}
Prove Proposition~\ref{ch31:prop-toolbox}(4): if $\Lap' = \Lap + w\,\mathbf{b}_e\mathbf{b}_e^\top$ with $w\ge0$ then $\lambda_i(\Lap')\ge\lambda_i(\Lap)$ for all $i$. Use the Courant--Fischer min--max formula. Then exhibit a graph where adding one edge strictly increases $\lambda_2$ and another where it leaves $\lambda_2$ unchanged.
\end{exercise}

\begin{exercise}[Harmonic = electrical]\label{ch31:ex-electrical}
Recompute the harmonic extension of Example~\ref{ch31:ex-harmonic} but ground vertex $4$ instead of vertex $1$ (set $x_4=0$, $x_5=1$). Interpret the result as voltages in a unit-resistance network, compute the energy $x^\top\Lap x$, and deduce the effective resistance between vertices $4$ and $5$.
\end{exercise}

\begin{exercise}[Normalized quadratic form]\label{ch31:ex-normalized}
Show that $x^\top\Lap_{\mathrm{sym}}x = \sum_{\{i,j\}\in E} w_{ij}\big(x_i/\sqrt{d_i} - x_j/\sqrt{d_j}\big)^2$. Conclude that $\Lap_{\mathrm{sym}}\succeq0$ and that its smallest eigenvalue is $0$ with eigenvector $\Adjm_D^{1/2}\mathbf{1}$. Why is this eigenvector, rather than $\mathbf{1}$, the ``constant'' of the normalized world?
\end{exercise}

\begin{exercise}[Spectral upper bound on the cut]\label{ch31:ex-spectral-cut}
Let $S\subset V$ with $\abs{S}=s$. Using $\mathbf{1}_S^\top\Lap\,\mathbf{1}_S = w(S,\bar S)$ and the Rayleigh quotient, prove $w(S,\bar S) \ge \lambda_2 \cdot \frac{s(n-s)}{n}$ for an unweighted graph. (Hint: project $\mathbf{1}_S$ onto $\mathbf{1}^\perp$ and bound its Rayleigh quotient below by $\lambda_2$.) This is the easy half of the spectral story; the hard half is Cheeger.
\end{exercise}

\begin{exercise}[Implement and time it]\label{ch31:ex-implement}
Implement \textsc{ComponentsByRank} (Algorithm~\ref{ch31:alg-rank}) and compare it to BFS on random graphs of increasing size. Plot running time versus $n$ and confirm the $\Th(n^3)$ versus $\Oh(n+m)$ gap. Then build $\Lap$ as a sparse matrix and time a single \textsc{HarmonicExtend} solve using a conjugate-gradient routine; observe how the iteration count grows with $\lambda_n/\lambda_2$.
\end{exercise}

\begin{exercise}[Trace the elimination]\label{ch31:ex-elim-trace}
Reproduce the Gaussian-elimination trace of the ``Trace it yourself'' box following Algorithm~\ref{ch31:alg-rank}: carry the running graph's $\Lap$ to upper-triangular form by hand, recording each pivot. Confirm the four pivots $2,\tfrac32,\tfrac73,\tfrac47$ and the single zero row. Then delete edge $\{3,5\}$ and re-run; show a second zero row appears and the pivot product drops, matching $k=2$.
\end{exercise}

\begin{exercise}[Spring energy minimization]\label{ch31:ex-springs}
Read Figure~\ref{ch31:fig-springs} literally. Pin beads $1$ and $3$ at heights $0$ and $1$ in a triangle $K_3$ with unit springs and let bead $2$ settle. Write the energy $\tfrac12 x^\top\Lap x$ as a function of $x_2$, minimize it, and check that the minimizer is the harmonic value $x_2=\tfrac12$. Relate the stored energy at the minimum to the effective resistance between the two pinned beads.
\end{exercise}
